\documentclass[11pt]{article}

\usepackage[utf8]{inputenc}
\usepackage[spanish]{babel}
\usepackage[margin=2.5cm,headheight=14pt]{geometry}
\usepackage{graphicx}
\usepackage{microtype}
\usepackage{booktabs}
\usepackage{longtable}
\usepackage{tikz}
\usepackage{listings}
\usepackage{xcolor}
\usepackage{hyperref}
\usepackage{fancyhdr}
\usepackage{titlesec}
\usepackage{enumitem}
\usetikzlibrary{shapes.geometric, arrows.meta, positioning}

\usepackage{array}
\newcolumntype{L}[1]{>{\raggedright\arraybackslash}p{#1}}
\emergencystretch=3em

\hypersetup{
    colorlinks=true,
    linkcolor=black,
    citecolor=black,
    urlcolor=blue,
    pdftitle={Proyecto 1 - Computacion Paralela y Distribuida},
}

\pagestyle{fancy}
\fancyhf{}
\fancyhead[L]{Screensaver de fluidos - OpenMP}
\fancyhead[R]{\thepage}
\renewcommand{\headrulewidth}{0.4pt}

\definecolor{termbg}{RGB}{20,20,20}
\definecolor{termfg}{RGB}{0,255,120}
\lstdefinestyle{terminal}{
    backgroundcolor=\color{termbg},
    basicstyle=\color{termfg}\ttfamily\small,
    breaklines=true,
    frame=single,
    framerule=0pt,
    xleftmargin=4pt,
    xrightmargin=4pt,
}
\lstdefinestyle{code}{
    basicstyle=\ttfamily\small,
    breaklines=true,
    frame=single,
    numbers=left,
    numberstyle=\tiny,
}

\begin{document}

% ============================================================
% CARATULA
% ============================================================
\begin{titlepage}
    \centering
    \vspace*{1cm}

    {\large Universidad del Valle de Guatemala\par}
    {\large Facultad de Ingenieria\par}
    {\large Departamento de Ciencias de la Computacion\par}
    {\large Computacion Paralela y Distribuida - Seccion 20\par}
    \vspace{0.3cm}

    \vspace{3cm}

    {\Huge\bfseries Proyecto \#1\par}
    \vspace{0.5cm}
    {\Large Screensaver de fluidos paralelizado con OpenMP\par}
    \vspace{0.3cm}
    {\large Simulacion de Navier-Stokes (metodo Stable Fluids) con n-cuerpos\par}

    \vspace{3cm}

    {\large\bfseries Integrantes\par}
    \vspace{0.3cm}
    {\large José Antonio Mérida Castejón\par}
    {\large Javier Eduardo España Pacheco\par}
    {\large Angel Esteban Esquit Hernández\par}

    \vfill

    {\large Guatemala, 2 de septiembre de 2026\par}
\end{titlepage}

\tableofcontents
\clearpage

% ============================================================
% INTRODUCCION
% ============================================================
\section{Introduccion}

Este informe documenta el desarrollo de un screensaver de fluidos,
paralelizado con OpenMP. El programa simula un fluido incompresible en 2D
usando el metodo Stable Fluids de Jos Stam, con fuentes de tinta que se
mueven por gravedad mutua (n-cuerpos).

El objetivo del proyecto es tomar una version secuencial funcional y
mejorarla de forma iterativa, midiendo speedup y eficiencia en cada paso.
Cada paso agrega una sola mejora (algoritmica o de paralelizacion) para que
su impacto se pueda medir por separado.

% ============================================================
% ANTECEDENTES
% ============================================================
\section{Antecedentes}

OpenMP es una API para programacion paralela de memoria compartida en
C/C++/Fortran. Funciona mediante directivas de compilador (\texttt{\#pragma
omp}) que le indican al compilador que partes del codigo se pueden repartir
entre varios hilos. A diferencia de MPI (memoria distribuida), todos los
hilos de OpenMP comparten el mismo espacio de memoria, lo cual simplifica
la programacion pero exige cuidado con condiciones de carrera.

El problema elegido, simulacion de fluidos, es un caso clasico de
paralelizacion de memoria compartida: la malla se puede dividir en
regiones, y (con el cuidado correcto de dependencias) cada region se puede
actualizar en un hilo distinto. El metodo de Stam (\emph{Stable Fluids})
resuelve las ecuaciones de Navier-Stokes con pasos incondicionalmente
estables, lo que permite pasos de tiempo grandes sin que la simulacion
explote numericamente \cite{stam1999}.

Para el movimiento de las fuentes de tinta se implemento un sistema de
n-cuerpos con gravedad mutua, acelerado con un quadtree Barnes-Hut
\cite{barneshut1986}, que reduce el costo de $O(n^2)$ a $O(n \log n)$.

% ============================================================
% CUERPO
% ============================================================
\section{Descripcion del problema}

El programa dibuja un lienzo de al menos 640x480 px donde varias fuentes
de tinta de colores inyectan color y momentum en un campo de velocidad.
El campo de velocidad se resuelve cada frame (difusion, adveccion,
proyeccion de Hodge) y arrastra la tinta, generando remolinos. Las fuentes
mismas se mueven por atraccion gravitacional mutua y rebotan en los bordes
de la malla.

Parametros configurables por linea de comandos: resolucion de la malla
(\texttt{-n}), cantidad de fuentes (\texttt{-f}), tamano de ventana
(\texttt{-W}/\texttt{-H}), semilla aleatoria (\texttt{-s}), viscosidad
(\texttt{-v}), difusion (\texttt{-d}) y pantalla completa (\texttt{-p}).
Los valores por defecto y rangos validos estan en la Tabla
\ref{tab:cli-params}.

\begin{table}[ht]
\centering
\begin{tabular}{lll}
\toprule
Flag & Rango & Default \\
\midrule
\texttt{-n} & 16 -- 1024 & 256 \\
\texttt{-f} & 1 -- 256 & 6 \\
\texttt{-W} & min. 640 & 1920 \\
\texttt{-H} & min. 480 & 1080 \\
\texttt{-v} & 0.0 -- 1.0 & 0.0 \\
\texttt{-d} & 0.0 -- 1.0 & 0.0001 \\
\bottomrule
\end{tabular}
\caption{Parametros de linea de comandos.}
\label{tab:cli-params}
\end{table}

\section{Metodologia: PCAM}

Se aplico el metodo PCAM (Particionamiento, Comunicacion, Aglomeracion,
Mapeo) para decidir que y como paralelizar:

\begin{description}[leftmargin=1.5cm]
    \item[Particionamiento] La malla de la simulacion se piensa como un
        conjunto de celdas independientes. Cada operacion por celda
        (inyeccion de fuente, disipacion, adveccion, proyeccion, render)
        es, en principio, una tarea separada.
    \item[Comunicacion] La mayoria de las operaciones solo leen el estado
        del frame anterior y escriben el nuevo, sin comunicacion entre
        celdas dentro del mismo paso. La excepcion es el solver de
        presion/difusion (Gauss-Seidel), que en su forma original lee
        valores ya actualizados de celdas vecinas en la misma iteracion,
        creando una dependencia estricta que impide paralelizar
        directamente.
    \item[Aglomeracion] En vez de repartir celda por celda (overhead de
        sincronizacion altisimo), se agrupan filas en bloques (\emph{chunks})
        que cada hilo procesa de corrido.
    \item[Mapeo] Los bloques se reparten entre los hilos disponibles via
        \texttt{\#pragma omp parallel for}, con la politica de reparto
        (\texttt{schedule}) elegida segun el tamano de la malla (ver
        Seccion \ref{sec:schedule}).
\end{description}

\section{Version secuencial}

La version secuencial (rama \texttt{sequential}) sirve de linea base fija
para medir speedup: una vez medida, no se vuelve a tocar, ni un commit
mas, para que el punto de comparacion no se mueva mientras se paraleliza.

Usa Gauss-Seidel fila por fila para el solver de presion/difusion y
fuerza bruta $O(f^2)$ para la gravedad del n-body (sin quadtree, cada
fuente contra cada otra fuente). El resto de la estructura del programa
(parseo de argumentos, asignacion de memoria, bucle principal, limpieza al
salir) es identica a las versiones paralelas — los pasos del roadmap solo
tocan el algoritmo del solver y el n-body, no la estructura general.

Antes de arrancar cualquier paralelizacion, esta version se probo:
compila limpio con \texttt{-Wall -Wextra -O2 -std=c11}, corre sin
crashear para el rango completo de \texttt{-n} y \texttt{-f}, y
\texttt{valgrind --leak-check=full} reporta 0 bytes perdidos y 0 errores
(la limpieza al salir libera los 12 campos de la simulacion, las fuentes
y las estrellas de fondo, ver \texttt{cleanup:} en \texttt{main.c}).
Tambien se probo el parseo de argumentos con entradas invalidas
(no numericas, fuera de rango, valores faltantes, flags desconocidos) —
ver Anexo 2 para el detalle de \texttt{read\_int}/\texttt{read\_float}.

\section{Paralelizacion iterativa}

Siguiendo el espiritu de OpenMP como herramienta iterativa, la
paralelizacion se hizo en pasos, cada uno en su propia rama, midiendo la
misma bateria de pruebas despues de cada uno.

\subsection{Paso 1: Barnes-Hut para n-body}

Se reemplazo el calculo de gravedad de fuerza bruta ($O(f^2)$) por un
quadtree Barnes-Hut ($O(f \log f)$). Este paso es puramente algoritmico,
sin OpenMP todavia. A las resoluciones de malla probadas, el solver de
fluidos domina tanto el costo que el cambio no se nota en el FPS total
(el n-body es una fraccion pequena del trabajo por frame).

\subsection{Paso 2: paralelizar los loops independientes}

Se agrego \texttt{\#pragma omp parallel for} a los sitios donde cada celda
de salida es independiente: inyeccion de fuentes, disipacion de tinta,
adveccion, proyeccion de Hodge y render. El solver de presion/difusion
seguia secuencial. Con 20 hilos, este paso ya da un salto importante de
FPS (Tabla \ref{tab:master}), pero el techo sigue siendo el solver.

\subsection{Paso 3: red-black Gauss-Seidel}
\label{sec:redblack}

El cambio mas grande de toda la cadena. El solver de Gauss-Seidel original
lee y escribe la misma malla en el mismo barrido, celda por celda en
orden, lo que crea una dependencia secuencial estricta. La solucion es un
recorrido \emph{red-black}: se divide la malla en un tablero de ajedrez
por paridad $(i+j) \bmod 2$. Los 4 vecinos de una celda de un color son
siempre del color opuesto, entonces todas las celdas del mismo color se
pueden actualizar en paralelo sin pisarse entre si. Cada iteracion pasa a
ser: actualizar todas las rojas (barrera implicita), actualizar todas las
negras (barrera implicita).

Este cambio finalmente paraleliza el hotspot principal del programa. El
resultado es el salto de rendimiento mas grande del proyecto: a la
configuracion default ($n=256$, $f=6$), el FPS pasa de 25.20 (Paso 2) a
88.27 (Paso 3), un $3.5\times$ adicional.

\subsection{Paso 4: ajuste de \texttt{schedule}}
\label{sec:schedule}

Con el hotspot ya paralelo, se probo si el \texttt{schedule(dynamic, 16)}
usado en todo el codigo era la mejor eleccion. Se cambiaron
temporalmente los \texttt{\#pragma} a \texttt{schedule(runtime)} (que se
puede configurar sin recompilar via la variable de entorno
\texttt{OMP\_SCHEDULE}) y se barrio \texttt{static}, \texttt{dynamic} y
\texttt{guided} con distintos tamanos de chunk, en malla chica y malla
grande (Tabla \ref{tab:schedule-sweep}).

El resultado principal: el \texttt{dynamic, 16} usado en todo el codigo es
buena eleccion en malla chica, pero es la peor opcion razonable en malla
grande. \texttt{static} (con el chunk por defecto del compilador) le saca
$2.5\times$ a \texttt{dynamic, 16} en $n=700$ (54.47 vs. 21.93 FPS).

Con este resultado, se aplico \texttt{schedule(static)} de verdad al
codigo de la rama \texttt{04-\allowbreak schedule-\allowbreak tuning},
reemplazando el \texttt{schedule(dynamic,\allowbreak\ 16)} hardcodeado en
todos los sitios del catalogo.

Al remedir la bateria estandar con el cambio real (no solo con
\texttt{OMP\_SCHEDULE}), el resultado es un \textbf{tradeoff}, no una
mejora universal. En malla chica \texttt{static} pierde un poco contra el
\texttt{dynamic, 16} anterior: la fila B6 baja de 88.27 a 80.68 FPS,
alrededor de un 8.6\% menos. En malla grande, en cambio, gana fuerte: la
fila I6 sube de 28.32 a 53.47 FPS (casi el doble) y la fila C6 sube de
45.79 a 68.57 FPS. Ver Tabla \ref{tab:master} para los numeros completos.

\subsection{Paso 5: ajuste de \texttt{collapse(2)} (cancelado)}

Se planeo comparar \texttt{collapse(2)} contra paralelizar solo el loop
externo, para ver el efecto en mallas chicas (menos filas que nucleos
disponibles). Este paso se cancelo antes de obtener numeros confiables y
no se incluye en los resultados de este informe.

\subsection{Paso 6: flags de compilacion}

Se probo cambiar el flag de optimizacion del compilador de \texttt{-O2}
(usado hasta ahora) a \texttt{-O3}, y adicionalmente \texttt{-march=native}
para autovectorizacion especifica del CPU de la maquina de pruebas. Los
resultados estan en la Tabla \ref{tab:o3}.

\texttt{-O3} da una ganancia real de $\sim$11\% en malla grande sobre
\texttt{-O2}, sin costo en malla chica. \texttt{-march=native} no aporta
nada adicional sobre \texttt{-O3} solo, lo que sugiere que los loops
calientes no se benefician de registros vectoriales mas anchos en este
CPU, o que el set de instrucciones base que asume GCC ya alcanza.

\section{Mecanismos de sincronizacion}

\begin{itemize}
    \item \textbf{\texttt{\#pragma omp parallel for}}: la forma principal
        de paralelismo en el proyecto. Reparte las iteraciones de un loop
        \texttt{for} entre los hilos disponibles, con una barrera
        implicita al final (todos los hilos esperan a que termine el
        ultimo antes de seguir).
    \item \textbf{\texttt{\#pragma omp parallel} + \texttt{\#pragma omp
        for}}: usado en el solver (Seccion \ref{sec:redblack}) para abrir
        una sola region paralela que envuelve las 20 iteraciones del
        Gauss-Seidel red-black, en vez de abrir/cerrar el team de hilos 40
        veces por frame (costo fijo de creacion de hilos evitado).
    \item \textbf{\texttt{\#pragma omp single}}: dentro de la region
        paralela del solver, aplica la condicion de frontera con un solo
        hilo mientras el resto espera en la barrera implicita del
        \texttt{single}, evitando que varios hilos escriban la misma
        celda de borde a la vez.
    \item \textbf{\texttt{schedule(dynamic, chunk)}}: politica de reparto
        de trabajo por defecto usada en todo el codigo. Reparte bloques
        de \texttt{chunk} iteraciones bajo demanda: un hilo que termina su
        bloque pide el siguiente en vez de quedarse ocioso esperando a
        los demas. La Seccion \ref{sec:schedule} muestra que no siempre es
        la mejor opcion.
    \item \textbf{\texttt{collapse(2)}}: en los loops anidados de
        adveccion, proyeccion y render, funde las dos dimensiones (fila,
        columna) en un solo rango repartible. Sin esto, una malla mas
        angosta que la cantidad de hilos dejaria nucleos sin trabajo.
    \item \textbf{Variables privadas (\texttt{private(...)})}: cada hilo
        necesita su propia copia de las variables escalares temporales
        del cuerpo del loop (por ejemplo, coordenadas de origen
        interpoladas en adveccion). Sin esto, los hilos pisarian el mismo
        valor compartido y corromperian el resultado.
\end{itemize}

No se usaron mutex ni secciones criticas explicitas: el diseno evita
condiciones de carrera por construccion (cada celda de salida se escribe
una sola vez, por un solo hilo, en cada paso), asi que las barreras
implicitas de \texttt{omp for} y el \texttt{single} para la frontera son
suficientes.

\section{Resultados}

\subsection{Tabla maestra}

La Tabla \ref{tab:master} muestra el FPS promedio por configuracion (malla
$\times$ fuentes), en cada paso de la cadena, con 20 hilos donde aplica.
Metodologia: cada fila se corrio con semilla fija (\texttt{-s 42}), se
promedio el FPS ignorando las primeras lineas de calentamiento.

\begin{table}[ht]
\centering
\small
\begin{tabular}{lrrrrr}
\toprule
Fila & \texttt{sequential} & \texttt{01-rb-tree} & \texttt{02-omp-loops} & \texttt{03-omp-solver} & \texttt{04-schedule} \\
& & & & (\texttt{dynamic,16}) & (\texttt{static}) \\
\midrule
A6 (n=64, f=6)    & 17.21 & 17.19 & 118.40 & 88.04 & 83.87 \\
A12 (n=64, f=12)  & 17.16 & 17.20 & 117.43 & 88.31 & 84.39 \\
B6 (n=256, f=6)   & 10.52 & 10.51 & 25.20  & 88.27 & 80.68 \\
B12 (n=256, f=12) & 10.95 & 10.95 & 25.24  & 88.66 & 80.92 \\
C6 (n=512, f=6)   & 3.75  & 3.76  & 6.72   & 45.79 & 68.57 \\
C12 (n=512, f=12) & 4.30  & 4.30  & 6.89   & 52.33 & 69.53 \\
H6 (n=600, f=6)   & 2.69  & 2.69  & 4.46   & 33.72 & 61.41 \\
H12 (n=600, f=12) & 3.09  & 3.09  & 4.79   & 41.39 & 62.29 \\
I6 (n=700, f=6)   & 1.90  & 1.91  & 3.00   & 28.32 & 53.47 \\
I12 (n=700, f=12) & 2.10  & 2.10  & 3.21   & 35.57 & 54.76 \\
D6 (n=1024, f=6)  & 0.92  & 0.92  & 1.15   & 5.10  & 4.64  \\
D12 (n=1024, f=12)& 1.00  & 0.99  & 1.23   & 6.32  & 3.21  \\
E (n=256, f=4)    & 7.52  & 7.52  & 24.39  & 87.09 & n/m \\
F (n=256, f=64)   & 11.39 & 11.39 & 25.36  & 87.28 & n/m \\
G (n=256, f=256)  & 11.37 & 11.38 & 25.30  & 85.07 & n/m \\
\bottomrule
\end{tabular}
\caption{FPS promedio por paso de la cadena de paralelizacion. n/m = no
medido con el schedule nuevo todavia.}
\label{tab:master}
\end{table}

En la configuracion default ($n=256$, $f=6$), el FPS pasa de 10.52
(secuencial) a 88.27 (Paso 3), un speedup acumulado de $\sim 8.4\times$.
El Paso 4 (\texttt{static}) es un tradeoff frente a Paso 3: pierde un
poco en malla chica pero gana fuerte en malla grande (hasta
$\sim +89\%$ en n=700, Tabla \ref{tab:master}).

\subsection{Escalabilidad por hilos}

La Tabla \ref{tab:threads} muestra el FPS variando \texttt{OMP\_NUM\_THREADS}
en la fila B (n=256, f=6). Antes del Paso 3 (solver todavia secuencial), el
escalado se estanca alrededor de $2.6\times$ (techo de Amdahl: el solver
secuencial limita cuanto puede acelerar el resto). Desde el Paso 3 en
adelante, con el solver ya paralelo, el escalado sube a $\sim 6\times$.

\begin{table}[ht]
\centering
\begin{tabular}{lrrr}
\toprule
Hilos & FPS \texttt{02-omp-loops} & FPS \texttt{03-omp-solver} & FPS \texttt{04-schedule} \\
\midrule
1  & 9.95  & 14.72 & 14.72 \\
2  & 14.94 & 25.39 & 25.52 \\
4  & 20.33 & 45.94 & 45.91 \\
8  & 23.57 & 67.84 & 68.27 \\
16 & 25.52 & 83.27 & 83.66 \\
20 & 25.57 & 87.76 & 87.72 \\
\bottomrule
\end{tabular}
\caption{FPS vs. cantidad de hilos, fila B (n=256, f=6).}
\label{tab:threads}
\end{table}

Tambien se corrio la misma curva en una malla grande (n=700, f=12),
comparando el \texttt{schedule} actual contra \texttt{static} (el mejor
encontrado para malla grande, Tabla \ref{tab:threads-700}). Con
\texttt{static}, 20 hilos da 48.11 FPS contra los 34.59 FPS del
\texttt{dynamic, 16} actual en la misma fila: +39\% solo cambiando el
schedule.

\begin{table}[ht]
\centering
\begin{tabular}{lrr}
\toprule
Hilos & FPS \texttt{static} & FPS \texttt{dynamic,16} (actual) \\
\midrule
1  & 6.82  & -- \\
2  & 15.25 & -- \\
4  & 30.58 & -- \\
8  & 41.02 & -- \\
16 & 53.80 & -- \\
20 & 48.11 & 34.59 \\
\bottomrule
\end{tabular}
\caption{FPS vs. hilos en malla grande (n=700, f=12).}
\label{tab:threads-700}
\end{table}

Con \texttt{static} el pico esta en 16 hilos, no en 20: a 700 filas, el
reparto estatico en 20 partes no cae tan parejo como en 16.

\subsection{Barrido de \texttt{schedule}/chunk}

\begin{table}[ht]
\centering
\begin{tabular}{lrr}
\toprule
\texttt{OMP\_SCHEDULE} & FPS n=256 & FPS n=700 \\
\midrule
\texttt{static} (default)          & 83.40 & \textbf{54.47} \\
\texttt{static,4}                  & 83.57 & 43.01 \\
\texttt{static,16}                 & 84.53 & 46.87 \\
\texttt{static,64}                 & 72.89 & 44.42 \\
\texttt{dynamic,1}                 & 31.19 & 5.55  \\
\texttt{dynamic,4}                 & 64.36 & 15.77 \\
\texttt{dynamic,16} (actual)       & 87.58 & 21.93 \\
\texttt{dynamic,64}                & 76.88 & 2.02  \\
\texttt{guided} (default)          & \textbf{96.80} & 49.42 \\
\texttt{guided,16}                 & 94.71 & 3.43  \\
\bottomrule
\end{tabular}
\caption{Barrido de schedule/chunk, 20 hilos, f=6.}
\label{tab:schedule-sweep}
\end{table}

\subsection{Flags de build}

\begin{table}[ht]
\centering
\begin{tabular}{lrr}
\toprule
Build & FPS n=256,f=6 & FPS n=700,f=12 \\
\midrule
\texttt{-O2} (actual)        & 83.74 & 32.18 \\
\texttt{-O3}                 & 83.16 & 35.72 \\
\texttt{-O3 -march=native}   & 84.03 & 35.76 \\
\bottomrule
\end{tabular}
\caption{FPS por flag de optimizacion, 20 hilos.}
\label{tab:o3}
\end{table}

% ============================================================
% CONCLUSIONES Y RECOMENDACIONES
% ============================================================
\section{Conclusiones}

\begin{itemize}
    \item Paralelizar el solver de presion/difusion (red-black
        Gauss-Seidel) fue, por lejos, el cambio de mayor impacto: FPS
        default subio de 25.20 a 88.27 ($3.5\times$), y la escalabilidad
        por hilos paso de $\sim 2.6\times$ a $\sim 6\times$.
    \item Antes de paralelizar el hotspot real, paralelizar el resto del
        programa (Paso 2) ayuda, pero el techo de Amdahl limita cuanto se
        puede ganar mientras el cuello de botella siga secuencial.
    \item La eleccion de \texttt{schedule} de OpenMP importa mas de lo
        esperado y no tiene un ganador universal: \texttt{dynamic, 16}
        era mejor en malla chica, \texttt{static} gana fuerte en malla
        grande ($2.5\times$ en el barrido, hasta $+89\%$ al aplicarlo de
        verdad al codigo). Se aplico \texttt{static} a la rama
        \texttt{04-schedule-tuning}; el resultado real confirma el
        tradeoff: pierde $\sim$8\% en malla chica para ganar hasta
        $\sim$89\% en malla grande.
    \item Cambiar \texttt{-O2} a \texttt{-O3} da una ganancia gratuita de
        $\sim$11\% en malla grande sin tocar codigo. \texttt{-march=native}
        no aporta nada extra sobre \texttt{-O3} en este CPU.
    \item Barnes-Hut para el n-body no se nota en el FPS a las
        resoluciones probadas porque el solver de fluidos domina el costo
        total; el algoritmo de n-body es una fraccion pequena del trabajo
        por frame.
\end{itemize}

\section{Recomendaciones}

\begin{itemize}
    \item Ya se aplico \texttt{static} en \texttt{04-schedule-tuning}. Si
        el uso real del programa favorece mallas chicas ($n \leq 256$),
        vale la pena reconsiderar, ya que ahi \texttt{dynamic, 16} sigue
        siendo mejor; para mallas grandes, \texttt{static} es la mejor
        eleccion medida. Una opcion futura es dejarlo configurable via
        \texttt{schedule(runtime)} y elegir en base a \texttt{-n} al
        arrancar.
    \item Cambiar \texttt{-O2} a \texttt{-O3} en el \texttt{Makefile}: es
        una ganancia gratuita sin riesgo.
    \item Retomar el Paso 5 (\texttt{collapse(2)}) con una metodologia mas
        controlada (por ejemplo, sin renderizado en pantalla completa) si
        se quiere ese numero para mallas chicas con pocos nucleos.
    \item Para mallas mayores a $n=1024$ o cantidades de fuentes mayores a
        256, medir de nuevo: el comportamiento del schedule optimo podria
        cambiar fuera del rango probado.
\end{itemize}

% ============================================================
% BIBLIOGRAFIA
% ============================================================
\begin{thebibliography}{9}

\bibitem{stam1999}
Stam, J. (1999). \emph{Stable Fluids}. En Proceedings of the 26th Annual
Conference on Computer Graphics and Interactive Techniques (SIGGRAPH '99),
121-128.

\bibitem{barneshut1986}
Barnes, J., \& Hut, P. (1986). \emph{A hierarchical $O(N \log N)$
force-calculation algorithm}. Nature, 324(6096), 446-449.

\bibitem{ompspec}
OpenMP Architecture Review Board. (2021). \emph{OpenMP Application
Programming Interface, Version 5.2}.
\url{https://www.openmp.org/specifications/}

\bibitem{chapman2007}
Chapman, B., Jost, G., \& Van Der Pas, R. (2007). \emph{Using OpenMP:
Portable Shared Memory Parallel Programming}. MIT Press.

\end{thebibliography}

% ============================================================
% ANEXO 1: DIAGRAMA DE FLUJO
% ============================================================
\clearpage
\appendix
\section{Anexo 1: Diagrama de flujo del programa}

\begin{center}
\resizebox{!}{0.83\textheight}{%
\begin{tikzpicture}[
    node distance=0.9cm and 1.5cm,
    box/.style={rectangle, draw, rounded corners, minimum width=4.2cm, minimum height=0.9cm, align=center, font=\small},
    decision/.style={diamond, draw, aspect=2.5, minimum width=3.6cm, align=center, font=\small},
    par/.style={rectangle, draw, rounded corners, minimum width=4.2cm, minimum height=0.9cm, align=center, font=\small, fill=blue!12},
    arr/.style={-{Stealth[length=2mm]}}
]

\node[box] (start) {Inicio};
\node[box, below=of start] (args) {Captura de argumentos\\(\texttt{parse\_arguments})};
\node[decision, below=of args] (argsok) {\textcolor{black}{args validos?}};
\node[box, right=of argsok] (argserr) {Imprime error /\\ayuda y termina};
\node[box, below=of argsok] (alloc) {Programacion defensiva:\\reserva de memoria\\(\texttt{allocate\_fields}, malloc)};
\node[decision, below=of alloc] (allocok) {memoria OK?};
\node[box, right=of allocok] (allocerr) {Libera lo reservado\\y termina};
\node[box, below=of allocok] (sdl) {Inicializa SDL,\\ventana, renderer};
\node[box, below=of sdl] (loop) {Bucle principal\\(por frame)};

\node[par, below=of loop] (nbody) {Mueve fuentes (n-body,\\Barnes-Hut)};
\node[par, below=of nbody] (inject) {Inyecta tinta/fuerza\\(\texttt{inject\_sources})};
\node[par, below=of inject] (vel) {Resuelve velocidad\\(seccion paralela:\\difusion + proyeccion,\\red-black Gauss-Seidel)};
\node[par, below=of vel] (advect) {Advecta tinta\\(\texttt{collapse(2)})};
\node[par, below=of advect] (dissip) {Disipa tinta};
\node[par, below=of dissip] (render) {Renderiza frame\\(\texttt{collapse(2)})};
\node[box, below=of render] (fps) {Mide y despliega FPS};
\node[decision, below=of fps] (quit) {ESC/Q o\\cerrar ventana?};
\node[box, right=of quit] (cleanup) {Libera memoria,\\destruye ventana,\\termina};

\draw[arr] (start) -- (args);
\draw[arr] (args) -- (argsok);
\draw[arr] (argsok) -- node[left]{no} (argserr);
\draw[arr] (argsok) -- node[right]{si} (alloc);
\draw[arr] (alloc) -- (allocok);
\draw[arr] (allocok) -- node[left]{no} (allocerr);
\draw[arr] (allocok) -- node[right]{si} (sdl);
\draw[arr] (sdl) -- (loop);
\draw[arr] (loop) -- (nbody);
\draw[arr] (nbody) -- (inject);
\draw[arr] (inject) -- (vel);
\draw[arr] (vel) -- (advect);
\draw[arr] (advect) -- (dissip);
\draw[arr] (dissip) -- (render);
\draw[arr] (render) -- (fps);
\draw[arr] (fps) -- (quit);
\draw[arr] (quit) -- node[right]{si} (cleanup);
\draw[arr] (quit.west) to[out=180,in=180] node[left]{no, siguiente frame} (loop.west);

\end{tikzpicture}%
}
\end{center}

Los bloques en azul son las secciones paralelizadas con OpenMP. La
sincronizacion entre ellas es implicita: cada \texttt{\#pragma omp for}
tiene una barrera al final, y la seccion de velocidad usa
\texttt{\#pragma omp single} para la condicion de frontera dentro de su
propia region paralela (ver Seccion \ref{sec:redblack}).

\subsection*{Que pasa dentro de cada bloque, un frame a la vez}

El diagrama de arriba muestra el orden, pero no el detalle de cada
bloque. Esta seccion entra a ese detalle, bloque por bloque, en el orden
real de ejecucion dentro de \texttt{main.c}.

\paragraph{1. Mueve fuentes (n-body, Barnes-Hut) — \texttt{update\_nbody\_sources}.}
Reconstruye un quadtree desde cero cada frame (las fuentes se mueven, no
tiene sentido mantenerlo entre frames). Primero inserta las $f$ fuentes
una por una (\texttt{bh\_insert}, recursivo, subdivide un nodo hoja en 4
cuando cae una segunda fuente adentro). Despues, para cada fuente,
recorre el arbol (\texttt{bh\_accumulate}): si un nodo esta lo bastante
lejos en relacion a su tamano ($\theta = 0.5$), se trata como un solo
punto de masa en vez de bajar a sus 4 hijos — eso es lo que baja el costo
de $O(f^2)$ a $O(f \log f)$. Con la aceleracion acumulada, integra
velocidad y posicion, y aplica un rebote amortiguado si una fuente
llegaria a salir del area valida de inyeccion. Esta etapa sigue
secuencial en toda la cadena: el arbol se construye una vez por frame y
depende del resultado del frame anterior en cada insercion, no vale la
pena paralelizarla a la cantidad de fuentes que se usan tipicamente
($f \leq 256$).

\paragraph{2. Inyecta tinta/fuerza — \texttt{inject\_sources}.}
Primero limpia los buffers "previous" de velocidad y tinta (van a
recibirse de nuevo este frame). Para cada fuente, calcula dos gaussianas
centradas en su posicion (continua, no la celda entera): una ancha que
da el color/empuje general, y una angosta ("core") que agrega un nucleo
brillante casi blanco en el centro. La direccion del chorro rota con
\texttt{cos(phase)}/\texttt{sin(phase)}, generando el efecto espiral. No
esta paralelizada: el radio de inyeccion por fuente es chico (unas pocas
celdas) y varias fuentes pueden solapar en la misma celda si estan cerca,
lo que forzaria una seccion critica o reduccion — no vale la pena al
tamano tipico de \texttt{-f}.

\paragraph{3. Resuelve velocidad — \texttt{velocity\_step}.}
El bloque mas caro del frame, y donde vive el paralelismo real. Es una
secuencia fija: fuerza $\to$ difusion $\to$ proyeccion $\to$
autoadveccion $\to$ proyeccion otra vez (difundir y autoadvectar cada uno
reintroduce algo de divergencia, por eso se proyecta dos veces). La
difusion y cada proyeccion llaman a \texttt{solve\_linear}, que abre
\emph{una sola} region \texttt{\#pragma omp parallel} y adentro corre las
20 iteraciones de Gauss-Seidel red-black: 20 veces (actualizar celdas
rojas, barrera implicita, actualizar celdas negras, barrera implicita).
Abrir la region paralela una sola vez en vez de 40 veces (20 iteraciones
$\times$ 2 colores) evita el costo fijo de crear el team de hilos
repetidamente, que en mallas chicas se nota. La condicion de frontera se
aplica con \texttt{\#pragma omp single} dentro de esa misma region: un
solo hilo la escribe mientras los demas esperan en la barrera implicita
del \texttt{single}, sin salir de la region paralela.

\paragraph{4. Advecta tinta — \texttt{ink\_step} (una vez por canal R, G, B).}
Adveccion semi-Lagrangiana: para cada celda destino, calcula de donde
vino ese fluido dando un paso hacia atras a lo largo de su propio campo
de velocidad, y agarra el valor de ahi por interpolacion bilineal de las
4 celdas reales alrededor de ese punto (fraccional). Cada celda destino
solo lee el campo anterior, nunca lo escribe, asi que es paralelizable
sin ningun cuidado extra. Usa \texttt{\#pragma omp parallel for
collapse(2)}: funde el loop de \texttt{i} y el de \texttt{j} en un solo
rango repartible, para que una malla angosta (menos filas que hilos)
siga repartiendo trabajo entre todos. Se llama 3 veces (una por canal de
color), cada una independiente de las otras.

\paragraph{5. Disipa tinta — \texttt{dissipate\_ink}.}
Multiplica cada celda de cada canal por \texttt{DISSIPATION} (0.995), sin
eso la tinta se acumularia hasta saturar la pantalla a blanco. Loop plano
sobre \texttt{celdas\_totales}, cada celda independiente,
\texttt{\#pragma omp parallel for}.

\paragraph{6. Renderiza frame — \texttt{render\_ink}.}
Recorre cada pixel de la textura (no cada celda de la malla — la
resolucion de ventana suele ser mayor que \texttt{-n}), interpola
bilinealmente la densidad de tinta en esa posicion, aplica un mapeo de
tono estilo Reinhard y un ajuste de saturacion, y calcula el canal alfa a
partir del brillo maximo de los 3 canales (para que el fondo se vea donde
no hay tinta). Tambien usa \texttt{collapse(2)}: sin el, con una ventana
mas angosta que alta (o viceversa) los hilos no se repartirian parejo
solo sobre el loop externo.

\paragraph{7. Mide FPS y actualiza pantalla.}
Cuenta frames acumulados y, cada $\sim$0.5s de tiempo real, calcula el
promedio, lo imprime a \texttt{stdout} y lo pone en el titulo de la
ventana. No es parte del calculo de la simulacion, es solo
instrumentacion.

% ============================================================
% ANEXO 2: CATALOGO DE FUNCIONES
% ============================================================
\clearpage
\section{Anexo 2: Catalogo de funciones}

\subsection*{\texttt{config.c}}
\begin{longtable}{L{4cm}L{4.1cm}L{2.2cm}L{3.7cm}}
\toprule
Funcion & Entradas & Salidas & Descripcion \\
\midrule
\texttt{parse\_arguments} & \texttt{argc} (int), \texttt{argv} (char**) & \texttt{config} (Config*, por referencia), retorna int (1 ok, 0 error, -1 ayuda) & Recorre \texttt{argv} y llena \texttt{config} con los valores validados. \\
\texttt{read\_int} & \texttt{text} (char*), rango min/max & \texttt{out} (int*), retorna 1/0 & Convierte texto a int validando formato y rango. \\
\texttt{read\_float} & \texttt{text} (char*), rango min/max & \texttt{out} (float*), retorna 1/0 & Igual que \texttt{read\_int} pero para float. \\
\texttt{print\_help} & \texttt{program\_name} (char*) & ninguna (imprime a stdout) & Imprime la pantalla de uso del programa. \\
\bottomrule
\end{longtable}

\subsection*{\texttt{fields.c}}
\begin{longtable}{L{4cm}L{4.1cm}L{2.2cm}L{3.7cm}}
\toprule
Funcion & Entradas & Salidas & Descripcion \\
\midrule
\texttt{allocate\_fields} & \texttt{fields} (FluidFields*), \texttt{resolution} (int) & retorna int (1 ok, 0 error) & Reserva memoria para los 12 campos de la simulacion (velocidad, tinta, presion, divergencia). \\
\texttt{free\_fields} & \texttt{fields} (FluidFields*) & ninguna & Libera los 12 campos y pone los punteros en NULL. \\
\texttt{clear\_fields} & \texttt{fields} (FluidFields*) & ninguna & Pone todos los campos en cero (usado al reiniciar con la tecla R). \\
\bottomrule
\end{longtable}

\subsection*{\texttt{sources.c}}
\begin{longtable}{L{4cm}L{4.1cm}L{2.2cm}L{3.7cm}}
\toprule
Funcion & Entradas & Salidas & Descripcion \\
\midrule
\texttt{init\_sources} & \texttt{sources} (InkSource*), \texttt{count}, \texttt{resolution} (int) & ninguna & Inicializa posicion, velocidad, masa, color y fuerza de cada fuente. \\
\texttt{inject\_sources} & \texttt{sources} (InkSource*), \texttt{count} (int), \texttt{fields} (FluidFields*), \texttt{aspect} (float) & ninguna (modifica \texttt{fields}) & Deposita tinta y momentum de cada fuente en un vecindario gaussiano. \\
\texttt{dissipate\_ink} & \texttt{fields} (FluidFields*) & ninguna & Reduce la tinta de cada celda por un factor fijo (paralelo con \texttt{omp parallel for}). \\
\bottomrule
\end{longtable}

\subsection*{\texttt{solver.c}}
\begin{longtable}{L{4cm}L{4.1cm}L{2.2cm}L{3.7cm}}
\toprule
Funcion & Entradas & Salidas & Descripcion \\
\midrule
\texttt{apply\_boundary} & \texttt{bnd\_type} (int), \texttt{field} (float*) & ninguna & Aplica condiciones de frontera (rebote en paredes). \\
\texttt{add\_source} & \texttt{dest}, \texttt{source} (float*), \texttt{dt} (float), \texttt{n} (int) & ninguna & Suma el termino fuente a cada celda (paralelo). \\
\texttt{relax\_color} & \texttt{field}, \texttt{field\_prev} (float*), \texttt{a}, \texttt{inv\_c} (float), \texttt{parity} (int) & ninguna & Actualiza las celdas de una paridad del tablero red-black (paralelo, dentro de region \texttt{omp parallel}). \\
\texttt{solve\_linear} & \texttt{bnd\_type} (int), \texttt{field}, \texttt{field\_prev} (float*), \texttt{a}, \texttt{c} (float) & ninguna & Corre las 20 iteraciones de Gauss-Seidel red-black dentro de una sola region paralela. \\
\texttt{diffuse} & \texttt{bnd\_type} (int), \texttt{field}, \texttt{field\_prev} (float*), \texttt{diff}, \texttt{dt} (float), \texttt{n} (int) & ninguna & Arma los coeficientes y llama a \texttt{solve\_linear} para difusion implicita. \\
\texttt{advect} & \texttt{bnd\_type} (int), \texttt{field}, \texttt{field\_prev} (float*), \texttt{vel\_x}, \texttt{vel\_y} (float*), \texttt{dt} (float) & ninguna & Adveccion semi-Lagrangiana (paralelo, \texttt{collapse(2)}). \\
\texttt{project} & \texttt{vel\_x}, \texttt{vel\_y}, \texttt{pressure}, \texttt{divergence} (float*) & ninguna & Proyeccion de Hodge: mide divergencia, resuelve presion, resta el gradiente de la velocidad. \\
\texttt{swap\_fields} & \texttt{field\_a}, \texttt{field\_b} (float**) & ninguna (intercambia punteros) & Intercambia los punteros de dos buffers de campo. \\
\texttt{ink\_step} & \texttt{ink}, \texttt{ink\_prev} (float**), \texttt{vel\_x}, \texttt{vel\_y} (float*), \texttt{diff}, \texttt{dt} (float), \texttt{n} (int) & ninguna & Un paso completo (difusion + adveccion) para un canal de tinta. \\
\texttt{velocity\_step} & \texttt{fields} (FluidFields*), \texttt{viscosity}, \texttt{dt} (float) & ninguna & Un paso completo de velocidad: fuerza, difusion, proyeccion, autoadveccion, proyeccion. \\
\bottomrule
\end{longtable}

\subsection*{\texttt{nbody.c}}
\begin{longtable}{L{4cm}L{4.1cm}L{2.2cm}L{3.7cm}}
\toprule
Funcion & Entradas & Salidas & Descripcion \\
\midrule
\texttt{bh\_new\_node} & \texttt{cx}, \texttt{cy}, \texttt{half} (float) & retorna int (indice del nodo) & Crea un nodo nuevo del quadtree Barnes-Hut. \\
\texttt{bh\_quadrant} & \texttt{n} (QuadNode*), \texttt{x}, \texttt{y} (float) & retorna int (0-3) & Determina en que cuadrante cae un punto. \\
\texttt{bh\_subdivide} & \texttt{idx} (int) & ninguna & Crea los 4 hijos de un nodo. \\
\texttt{bh\_insert} & \texttt{idx}, \texttt{body} (int), \texttt{sources} (InkSource*), \texttt{depth} (int) & ninguna & Inserta una fuente en el arbol recursivamente. \\
\texttt{bh\_accumulate} & \texttt{idx}, \texttt{body} (int), \texttt{sources} (InkSource*), \texttt{acc\_x}, \texttt{acc\_y} (float*) & ninguna (acumula en \texttt{acc\_x}\allowbreak/\texttt{acc\_y}) & Recorre el arbol acumulando aceleracion gravitacional sobre una fuente. \\
\texttt{update\_nbody\_sources} & \texttt{sources} (InkSource*), \texttt{count}, \texttt{resolution} (int) & ninguna & Reconstruye el arbol, calcula gravedad, integra posicion/velocidad y aplica rebote en bordes. \\
\bottomrule
\end{longtable}

\subsection*{\texttt{render.c}}
\begin{longtable}{L{4cm}L{4.1cm}L{2.2cm}L{3.7cm}}
\toprule
Funcion & Entradas & Salidas & Descripcion \\
\midrule
\texttt{sample\_bilinear} & \texttt{field} (float*), \texttt{fx}, \texttt{fy} (float) & retorna float & Interpola bilinealmente un campo en una posicion continua. \\
\texttt{render\_ink} & \texttt{texture} (SDL\_Texture*), \texttt{fields} (FluidFields*), \texttt{width}, \texttt{height} (int) & ninguna & Vuelca la densidad de tinta a la textura SDL, un pixel a la vez (paralelo, \texttt{collapse(2)}). \\
\bottomrule
\end{longtable}

\subsection*{\texttt{background.c}}
\begin{longtable}{L{4cm}L{4.1cm}L{2.2cm}L{3.7cm}}
\toprule
Funcion & Entradas & Salidas & Descripcion \\
\midrule
\texttt{init\_background} & \texttt{stars} (Star*), \texttt{count}, \texttt{width}, \texttt{height} (int) & ninguna & Coloca estrellas de fondo en posiciones aleatorias. \\
\texttt{update\_background} & \texttt{stars} (Star*), \texttt{count} (int) & ninguna & Avanza el parpadeo de cada estrella. \\
\texttt{draw\_background} & \texttt{renderer} (SDL\_Renderer*), \texttt{stars} (Star*), \texttt{count} (int) & ninguna & Dibuja las estrellas de fondo. \\
\bottomrule
\end{longtable}

\subsection*{\texttt{utils.c}}
\begin{longtable}{L{4cm}L{4.1cm}L{2.2cm}L{3.7cm}}
\toprule
Funcion & Entradas & Salidas & Descripcion \\
\midrule
\texttt{random\_range} & \texttt{min}, \texttt{max} (float) & retorna float & Numero pseudoaleatorio uniforme en $[min, max]$. \\
\texttt{clamp} & \texttt{value}, \texttt{min}, \texttt{max} (float) & retorna float & Limita un valor al rango $[min, max]$. \\
\bottomrule
\end{longtable}

\subsection*{\texttt{main.c}}
\begin{longtable}{L{4cm}L{4.1cm}L{2.2cm}L{3.7cm}}
\toprule
Funcion & Entradas & Salidas & Descripcion \\
\midrule
\texttt{main} & \texttt{argc} (int), \texttt{argv} (char**) & retorna int (codigo de salida) & Punto de entrada: parsea argumentos, inicializa SDL y memoria, corre el bucle principal, libera todo al salir. \\
\bottomrule
\end{longtable}

% ============================================================
% ANEXO 3: BITACORA DE PRUEBAS
% ============================================================
\clearpage
\section{Anexo 3: Bitacora de pruebas}

\subsection*{Metodologia}

Cada configuracion se corrio con semilla fija (\texttt{-s 42}), ignorando
las primeras lineas de FPS (calentamiento/cache fria), y promediando el
resto. La bateria estandar (15 configuraciones, Tabla \ref{tab:matrix})
se repitio identica despues de cada paso de la cadena; son 15 mediciones
por paso, mas de las 10 minimas requeridas.

\begin{table}[ht]
\centering
\small
\begin{tabular}{lrrl}
\toprule
Fila & \texttt{-n} & \texttt{-f} & Motivo \\
\midrule
A6/A12   & 64   & 6/12  & malla chica, referencia rapida \\
B6/B12   & 256  & 6/12  & configuracion default \\
C6/C12   & 512  & 6/12  & malla grande, domina el solver \\
H6/H12   & 600  & 6/12  & rango de FPS medio observado \\
I6/I12   & 700  & 6/12  & extremo superior de ese rango \\
D6/D12   & 1024 & 6/12  & malla muy grande, limite superior \\
E        & 256  & 4     & pocas fuentes \\
F        & 256  & 64    & muchas fuentes \\
G        & 256  & 256   & limite superior de fuentes \\
\bottomrule
\end{tabular}
\caption{Bateria estandar de pruebas (15 configuraciones).}
\label{tab:matrix}
\end{table}

Para el barrido de \texttt{schedule} (Tabla \ref{tab:schedule-sweep}) se
probaron 10 combinaciones de politica/chunk, cada una en 2 tamanos de
malla: 20 mediciones. Para la escalabilidad por hilos (Tablas
\ref{tab:threads} y \ref{tab:threads-700}) se probaron 6 cantidades de
hilos por curva.

\subsection*{Captura de una corrida}

A continuacion, salida real de consola de una corrida de referencia
($n=256$, $f=6$, semilla 42), mostrando el arranque, calentamiento y
estabilizacion del FPS:

\begin{lstlisting}[style=terminal]
Fluid simulation (Navier-Stokes) - SEQUENTIAL version
  Grid         : 256 x 256 cells (65536 interior cells)
  Sources      : 6
  Window       : 1920 x 1080 px
  Seed         : 42
  visc / diff  : 0.00000 / 0.00010
FPS= 12.73
FPS= 11.41
FPS= 12.96
FPS= 13.80
FPS= 75.53
FPS= 83.46
FPS= 83.82
FPS= 82.31
FPS= 84.26
FPS= 85.41
FPS= 84.09
FPS= 63.32
FPS= 74.58
FPS= 70.27
FPS= 77.99
\end{lstlisting}

Nota: la etiqueta "SEQUENTIAL version" en el encabezado es texto residual
del archivo \texttt{main.c} (no se actualizo al paralelizar); el binario
usado en esta corrida es la version paralela con OpenMP habilitado (20
hilos).

\subsection*{Comando usado}

\begin{lstlisting}[style=code]
timeout 30s ./ss -n <N> -f <F> -s 42 2>/dev/null \
  | grep "FPS=" | tail -n +6 \
  | awk -F'= ' '{s+=$2; c++} END {printf "FPS promedio: %.2f\n", s/c}'
\end{lstlisting}

El resto de las mediciones completas (las 15 filas de la bateria estandar
por cada paso, el barrido de schedule y las curvas de escalabilidad) estan
documentadas en \texttt{docs/Roadmap.md} dentro del repositorio del
proyecto, junto con el historial de commits que muestra cuando se corrio
cada bateria.

\end{document}
